\documentclass[SE,lsstdraft,authoryear,toc]{lsstdoc}
\input{meta}

% Package imports go here.

% Local commands go here.

%If you want glossaries
%\input{aglossary.tex}
%\makeglossaries

\title{Investigation of Ellipticity as a Function of Degrees of Freedom on the Vera C. Rubin Observatory}

% This can write metadata into the PDF.
% Update keywords and author information as necessary.
\hypersetup{
    pdftitle={Ellipticity_DOFs},
    pdfauthor={RebekahPolen},
    pdfkeywords={}
}

\input{authors}

\setDocRef{SOTN-003}
\setDocUpstreamLocation{\url{https://github.com/lsst-so/sotn-003}}

\date{\vcsDate}

\setDocAbstract{%

This technote outlines the investigation into a potential relationship between ellipticity across the focal plane of the LSST Camera and the degrees of freedom of the Simonyi Survey Telescope. 
The aim of this investigation was ideally to find a relationship between the behavior of the ellipticity and a particular degree of freedom or combination of degrees of freedom. 
Such a finding would allow us to manipulate the optical system in order to correct for ellipticity when the behavior is contributing to the degradation of the point spread function. 
The results of this investigation reveal that, likely due to the highly nonlinear nature of second moments and ellipticity, any relationship to the degrees of freedom cannot be established. 
Rather, the investigation reveals that the common assumption that the degrees of freedom `start' near zero is usually false. 
Without knowing that initial conditions of the optical system, we cannot infer any potential relationship to the ellipticity behavior. 
The modeled ellipticity value integrated into Rubin TV rarely agrees with the measured ellipticity value, probably also due to the discrepancy of the initial conditions assumption. 

}

% Change history defined here.
% Order: oldest first.
% Fields: VERSION, DATE, DESCRIPTION, OWNER NAME.
% See LPM-51 for version number policy.
\setDocChangeRecord{%
  \addtohist{1}{YYYY-MM-DD}{Unreleased.}{Polen}
}


\begin{document}

% Create the title page.
\maketitle
% Frequently for a technote we do not want a title page  uncomment this to remove the title page and changelog.
% use \mkshorttitle to remove the extra pages

\section{Introduction}

The Vera C. Rubin Observatory Active Optics System (AOS) relies on curvature wavefront sensing to measure and correct for offsets on the Simonyi Survey Telescope. 
The AOS does this by manipulating the 50 degrees of freedom on the optical components of the telescope: the LSST Camera, the M2 secondary mirror, and the M1M3 mirror monolith consisting of the primary and tertiary mirrors. 
By ensuring that the optical components of the telescope are optimally aligned, we can improve image quality and reduce optical contributions to the point spread function. 

%However, 

\appendix

\section{Acknowledgements}

This material is based upon work supported in part by the National Science Foundation through Cooperative Agreements AST-1258333 and AST-2241526 and Cooperative Support Agreements AST-1202910 and AST-2211468 managed by the Association of Universities for Research in Astronomy (AURA), and the Department of Energy under Contract No.\ DE-AC02-76SF00515 with the SLAC National Accelerator Laboratory managed by Stanford University.
Additional Rubin Observatory funding comes from private donations, grants to universities, and in-kind support from LSST-DA Institutional Members.

% Include all the relevant bib files.
% https://lsst-texmf.lsst.io/lsstdoc.html#bibliographies
\section{References} \label{sec:bib}
\renewcommand{\refname}{} % Suppress default Bibliography section
\bibliography{local,lsst,lsst-dm,refs_ads,refs,books}

% Make sure lsst-texmf/bin/generateAcronyms.py is in your path
\section{Acronyms} \label{sec:acronyms}
\input{acronyms.tex}
% If you want glossary uncomment below -- comment out the two lines above
%\printglossaries





\end{document}
